\documentclass[10pt,a4paper]{amsart}
\usepackage[margin=19mm]{geometry}
\usepackage{amsmath,amssymb,mathtools}
\usepackage[scr=rsfs,cal=euler]{mathalfa}
\usepackage{xfrac}
\usepackage[unicode,hidelinks,bookmarks=false]{hyperref}
\newcommand{\ie}{i.e.\ }
\newcommand{\cf}{cf.\ }
\newcommand{\resp}{respectively\ }
\newcommand{\Subsection}{\subsection}
\providecommand{\nobreakdash}{\nobreak\hbox{-}}
\setlength{\emergencystretch}{2em}
\multlinegap0pt
\usepackage{xcolor}
\newtheorem{lem}{Lemma}
\title{Geodesic extensions of mechanical systems with nonholonomic constraints}
\author{Malika Belrhazi, Tom Mestdag}
\date{Source: arXiv:2404.11997v2 (CC BY 4.0)}
\begin{document}
\maketitle

\noindent\emph{Source and licence.} Geodesic extensions of mechanical systems with nonholonomic constraints. Version arXiv:2404.11997v2 (CC BY 4.0), distributed by arXiv under the Creative Commons Attribution 4.0 International licence, http://creativecommons.org/licenses/by/4.0/, as recorded in the arXiv licence metadata for this exact version and shown on the abstract page https://arxiv.org/abs/2404.11997v2. Full licence notice: \texttt{licenses/arxiv-2404.11997-CC-BY-4.0.txt}.\par\smallskip

\noindent\textit{Editorial scope.} Counted unit: the article's lem hulplemma (source0.tex lines 898--903). The article's complete original proof of it is delivered with it (source0.tex lines 905--973, one \texttt{proof} environment of 69 lines). No mathematics is written, repaired or re-derived: the statement and the proof are the article's own lines, in the article's own notation, and no retained line is substituted or re-flowed.  No further source-local context is needed: every cross-reference of every retained line resolves inside the two delivered spans.  Nothing was omitted from any retained span, so the package carries no omission record.  No retained line contains a citation, so the wrapper carries no bibliography and no bibliography entry.  No retained line contains a figure environment, an image-inclusion command or drawing code, so no image is delivered and the package declares no asset dependency.  Editorial formatting only, in addition: an AMS \texttt{amsart} preamble that restates the article's own declarations and macros used by the retained lines, namely the environments \texttt{lem} on separate counters, together with the standard packages those lines need.  The retained environments print in this document as Lemma 1 in order of appearance, whereas the article prints the same items with the numbering of its own full document.  The complete native source of the article as distributed in the arXiv e-print for this exact version is bundled unchanged as \texttt{sources/157851/source0.tex}.
\begin{lem} \label{hulplemma} At $0_q\in T_qQ$, we have 
$E_a^b(0)=\delta_a^b$, $E_a^i(0)=0$,
$$
\frac{\partial E^d_a}{\partial v^b}(0)=\frac{1}{2}R^d_{ab}(q)-\frac{1}{2}(\Gamma_{ab}^d(q)+\Gamma^d_{ba}(q)) \quad\text{and}\quad \frac{\partial E^i_a}{\partial v^b}(0)=\frac{1}{2}R^i_{ab}(q).
$$
\end{lem}

\begin{proof} The first two properties follow from the fact that, when $v_q=0_q$,   $T_{0_q}(\exp^{nh}_q)$ can be regarded as the inclusion.
We will compute the other expressions in two steps. 

Let $v_q = v^aX_a(q)\in T_qQ$. Consider a nonholonomic trajectory $c_{v_q}(t)$ and its derivative $\dot{c}_{v_q}(t)=V^a(t) X_a\left(c_{v_q}(t)\right) \in \mathcal{D}_{c_{v_q}(t)}$. The functions $V^a(t)$ then satisfy $\dot{V}^a(t)=-\Gamma_{b c}^a\left(c_{v_q}(t)\right) V^b(t) V^c(t)$.
Since in the purely kinetic case {$c_{v_q}(t)=\exp^{nh} \left(t v_q\right)$}, we may also write this as 
\[
\dot{c}_{v_q}(t)=E_a^\alpha(tv) v^a X_\alpha\left(c_{v_q}(t)\right).
\]
We get from this  
$$
E_{a}^d(t v_q) v^a=V^d(t) \quad \text{and} \quad E_a^i(t v_q) v^a=0.
$$
By applying $\frac{d}{dt}|_{t=0}$ on both sides of both equations we get 
$$
\frac{\partial E_a^d}{\partial v^\mu}(0) v^\mu v^a= \dot{V}^d(0)=-\Gamma^d_{ef}(q) v^e v^f
$$ and after symmetrization
$$ \frac{\partial E_a^d}{\partial v^c}(0)+\frac{\partial E_c^d}{\partial v^a}(0) =-(\Gamma_{ac}^d(q)+\Gamma_{ca}^d(q)).
$$
Likewise, it follows from 
$
{\partial E_a^i}/{\partial v^\mu}(0) v^\mu v^a = {\partial E_a^i}/{\partial v^b}(0) v^b v^a=0
$ that 
${\partial E_a^i}/{\partial v^b}(0)+{\partial E_b^i}/{\partial v^a}(0)=0$.


Now we compute the expression for the differences $ {\partial E_a^d}/{\partial v^c}(0)-{\partial E_c^d}/{\partial v^a}(0) $ and ${\partial E_a^i}/{\partial v^b}(0)-{\partial E_b^i}/{\partial v^a}(0)$. 
We know that in natural coordinates and in quasi-velocities, we get (respectively) 
\[
    T_{v_q}(\exp^{nh}_q)(u_q) = u^a\frac{\partial \epsilon^\alpha}{\partial v^a}(v) \frac{\partial}{\partial q^\alpha}|_{\tilde{q}} = E_a^\beta(v) u^a X_\beta(\tilde{q}).
		\]
Since $\{{\partial}/{\partial q^\alpha}|_{\tilde q}\}$ is a basis for $T_{\tilde q} Q$ there  exists an invertible matrix  $(A_\beta^\alpha(\tilde{q}))$ of functions such that
$X_\beta(\tilde{q})=A_\beta^\alpha(\tilde{q}) {\partial}/{\partial q^\alpha}|_{\tilde{q}}$. Therefore 
$ {\partial \epsilon^\alpha}/{\partial v^a}(v)=E_a^\beta(v) A_\beta^\alpha(\epsilon(v))$, and  in particular ${\partial \epsilon^\alpha}/{\partial v^a}(0)=E_a^\beta(0) A_\beta^\alpha(q)=A_a^\alpha(q)$.
When we take a $\partial/\partial v^b$-derivative on both sides of the first relation, we get
\[
    \frac{\partial^2 \epsilon^\alpha}{\partial v^a\partial v^b}(v)=\frac{\partial E_a^\beta}{\partial v^b}(v) A_\beta^\alpha(\epsilon(v))+E_a^\beta(v) \frac{\partial A_\beta^\alpha}{\partial q^\mu}(\epsilon(v))\frac{\partial \epsilon^\mu}{\partial v^b}(v),
	\]
and at  $v=0$:
\[
\frac{\partial^2 \epsilon^\alpha}{\partial v^a\partial v^b}(0)=\frac{\partial E_a^\beta}{\partial v^b}(0) A_\beta^\alpha(q)+\delta_a^\beta \frac{\partial A_\beta^\alpha}{\partial q^\mu}(q)
A^\mu_b(q).
\]
If we change in this last equality  $a$ and $b$, and subtract we get the identity
$$
 0=\left(\frac{\partial E_a^\beta}{\partial v^b}(0)-\frac{\partial E_b^\beta}{\partial v^a}(0)\right) A_\beta^\alpha(q)+\left( \frac{\partial A_a^\alpha}{\partial q^\mu}(q)A^\mu_b(q)- \frac{\partial A_b^\alpha}{\partial q^\mu}(q)A^\mu_a(q) \right).
$$
The last term is, in fact, related to the Lie bracket
$ R^\beta_{ba} X_\beta= [X_b,X_a]=(A^\mu_b{\partial A_a^\alpha}/{\partial q^\mu}- A^\mu_a{\partial A_b^\alpha}/{\partial q^\mu}) {\partial}/{\partial q^\alpha}$. 
The difference now becomes
$$ 
0=\left[\frac{\partial E_a^\beta}{\partial v^b}(0)-\frac{\partial E_b^\beta}{\partial v^a}(0) -R_{ab}^\beta(q)\right]A^\alpha_{\beta}(q).
$$
For $\beta=d$ and $\beta=i$, we get, respectively,
 $$
\frac{\partial E_a^d}{\partial v^b}(0)-\frac{\partial E_b^d}{\partial v^a}(0)=R_{ab}^d(q) ,
\qquad \frac{\partial E_a^i}{\partial v^b}(0)-\frac{\partial E_b^i}{\partial v^a}(0) =R_{ab}^i(q).
$$
We conclude that 
\begin{equation}
    \begin{cases}
\displaystyle      \frac{\partial E_a^d}{\partial v^b}(0)-\frac{\partial E_b^d}{\partial v^a}(0)=R_{ab}^d(q),\\[2mm]
   \displaystyle   \frac{\partial E_a^d}{\partial v^c}(0)+\frac{\partial E_c^d}{\partial v^a}(0) =-(\Gamma_{ac}^d(q)+\Gamma_{ca}^d(q)),
    \end{cases}\nonumber  \begin{cases}
\displaystyle        \frac{\partial E_a^i}{\partial v^b}(0)-\frac{\partial E_b^i}{\partial v^a}(0) =R_{ab}^i(q), \\[2mm]
\displaystyle        \frac{\partial E_a^i}{\partial v^b}(0)+\frac{\partial E_b^i}{\partial v^a}(0)=0 .
    \end{cases}
\end{equation}
This leads to the expressions in the statement of the Lemma.
\end{proof}

\end{document}
